\chapter{Contexte et problématique}

\section{Contexte du projet}
Le suivi des coopératives et de leurs bénéficiaires au sein de l'Axe Économie Sociale et
Solidaire repose aujourd'hui, selon les informations recueillies, en partie sur des
fichiers bureautiques (Excel, CSV) et des remontées manuelles, ce qui limite la
visibilité consolidée et complique la production de rapports fiables. \TODO{Confirmer
avec la Direction concernée l'état exact de l'existant avant publication du rapport
final — ne pas présenter cette description comme un fait officiel non vérifié.}

\section{Problématique}
Comment concevoir, avant tout développement, une spécification complète, cohérente et
sécurisée de la future plateforme de gestion des bénéficiaires — couvrant les besoins
fonctionnels, les contraintes techniques et les exigences de cybersécurité — à partir
d'un entretien structuré avec les parties prenantes, plutôt que d'un cahier des charges
générique ?

\section{Besoin métier}
La Fondation OCP a besoin d'un cadrage rigoureux, traçable et réutilisable des exigences
du futur système, garantissant que chaque exigence documentée provient effectivement
d'une réponse de partie prenante, d'une règle métier explicite, ou d'une hypothèse
clairement identifiée comme telle — et non d'une invention.

\section{Objectifs du projet}
\begin{itemize}
    \item Fournir un outil d'entretien structuré (IA ou questionnaire classique) dédié à
    ce projet précis, et non un chatbot généraliste.
    \item Générer un Cahier des Charges concis (3 à 5~pages) directement dérivé des
    réponses collectées.
    \item Produire une conception applicative de haut niveau et un diagramme de
    conception MVP unique.
    \item Générer des exigences de cybersécurité proportionnées aux risques
    effectivement identifiés durant l'entretien.
    \item Garantir qu'aucune information confidentielle de la Fondation OCP n'est
    inventée par le système.
\end{itemize}

\section{Périmètre}
Le périmètre de ce stage couvre la conception et le développement de
\textbf{Secure Requirements Studio}, l'outil d'ingénierie des exigences. Il ne couvre
\textbf{pas} le développement de la future Plateforme de Gestion des Bénéficiaires
elle-même, qui reste un projet ultérieur distinct, à charge de l'équipe de
développement qui reprendra les livrables produits par SRS.
